\originalchapter{8.05 作业10：Bell相关与角动量代数}
来源为Fall 2013 Assignment 10, Barton Zwiebach. 第1、2、5题有完整题面, 第3题所指的七道练习在MIT自己的Chapter 9原件可实际取得, 本篇逐道纳入. 解答均为\textbf{AI独立解答}. 第4题仅为Griffiths 4.57、p.197引用, 不补造. 原件公式竖线用页图核实, 不依赖失去绝对值的文本提取.
\originalsection{第1题：局域预定结果与Bell界}
\begin{exercise}
定义$C(a,b)=4\langle S_{1a}S_{2b}\rangle/\hbar^2$. 对单态完美反关联的局域预定自旋结果, 取$g=-4S_{1a}S_{1b}(1-4S_{1b}S_{1c}/\hbar^2)/\hbar^2$. (a)求$\langle g\rangle$. (b)由$|g|$证$|C(a,b)-C(a,c)|-C(b,c)\le1$. (c)量子单态求$C$. (d)三共面方向夹角$\theta_{ab}=\theta_{bc}=\theta,\theta_{ac}=2\theta$, 检验违反范围.
\end{exercise}
\begin{solution}
令$s_n=2S_{1n}/\hbar=\pm1$, 完美反关联给第二粒子结果$-s_n$, 所以$C(a,b)=-\langle s_as_b\rangle$. $g=-s_as_b+s_as_c$, 均值$C(a,b)-C(a,c)$. 因$1-s_bs_c$是0或2,$|g|=1-s_bs_c$, 均值$1+C(b,c)$. 三角不等式$|\langle g\rangle|\le\langle|g|\rangle$立即给题中Bell界. 对不同方向同时赋值与局域性是推导假设, 不能把不对易量子算符乘积当作同一组经典随机结果直接代进去.

量子单态$(\ket{+-}-\ket{-+})/\sqrt2$有$\langle\sigma_i\otimes\sigma_j\rangle=-\delta_{ij}$, 可对三个轴直接乘Pauli矩阵核, 交叉轴0. 故$C=-a\cdot b=-\cos\theta_{ab}$. 在$0\le\theta\le\pi/2$, Bell左侧为$|\cos2\theta-\cos\theta|+\cos\theta=1+2\cos\theta-2\cos^2\theta$. 因而$0<\theta<\pi/2$严格大于1, 两端恰等号; $\theta=\pi/3$达到3/2. 原件末句把端点也称违反不够精确, 此处明确修正. 绝对值在原PDF存在, 文本提取遗漏它会导致错误的“无违反”判断.
\begin{center}\begin{tikzpicture}
\begin{axis}[width=.82\textwidth,height=4.2cm,xlabel={$\theta$},ylabel={$F(\theta)$},grid=major,legend style={font=\small},legend pos=south east]
\addplot+[mark=none] coordinates{(0,1) (0.00654498469,1.00004284) (0.0130899694,1.00017133) (0.0196349541,1.00038544) (0.0261799388,1.00068512) (0.0327249235,1.00107025) (0.0392699082,1.00154074) (0.0458148929,1.00209644) (0.0523598776,1.00273717) (0.0589048623,1.00346276) (0.0654498469,1.00427299) (0.0719948316,1.0051676) (0.0785398163,1.00614633) (0.085084801,1.00720888) (0.0916297857,1.00835495) (0.0981747704,1.00958417) (0.104719755,1.01089619) (0.11126474,1.0122906) (0.117809725,1.01376699) (0.124354709,1.01532491) (0.130899694,1.0169639) (0.137444679,1.01868344) (0.143989663,1.02048304) (0.150534648,1.02236214) (0.157079633,1.02432016) (0.163624617,1.02635653) (0.170169602,1.02847063) (0.176714587,1.0306618) (0.183259571,1.03292939) (0.189804556,1.0352727) (0.196349541,1.03769103) (0.202894526,1.04018363) (0.20943951,1.04274974) (0.215984495,1.04538859) (0.22252948,1.04809936) (0.229074464,1.05088122) (0.235619449,1.05373332) (0.242164434,1.05665478) (0.248709418,1.05964471) (0.255254403,1.06270218) (0.261799388,1.06582625) (0.268344372,1.06901595) (0.274889357,1.07227031) (0.281434342,1.0755883) (0.287979327,1.0789689) (0.294524311,1.08241106) (0.301069296,1.0859137) (0.307614281,1.08947573) (0.314159265,1.09309604) (0.32070425,1.09677349) (0.327249235,1.10050692) (0.333794219,1.10429516) (0.340339204,1.10813702) (0.346884189,1.11203128) (0.353429174,1.11597671) (0.359974158,1.11997205) (0.366519143,1.12401603) (0.373064128,1.12810736) (0.379609112,1.13224473) (0.386154097,1.13642682) (0.392699082,1.14065228) (0.399244066,1.14491975) (0.405789051,1.14922784) (0.412334036,1.15357517) (0.41887902,1.15796031) (0.425424005,1.16238183) (0.43196899,1.1668383) (0.438513975,1.17132824) (0.445058959,1.17585018) (0.451603944,1.18040262) (0.458148929,1.18498405) (0.464693913,1.18959296) (0.471238898,1.1942278) (0.477783883,1.19888701) (0.484328867,1.20356904) (0.490873852,1.2082723) (0.497418837,1.21299519) (0.503963822,1.21773611) (0.510508806,1.22249345) (0.517053791,1.22726556) (0.523598776,1.23205081) (0.53014376,1.23684753) (0.536688745,1.24165406) (0.54323373,1.24646873) (0.549778714,1.25128983) (0.556323699,1.25611567) (0.562868684,1.26094455) (0.569413668,1.26577473) (0.575958653,1.27060449) (0.582503638,1.2754321) (0.589048623,1.28025579) (0.595593607,1.28507382) (0.602138592,1.28988443) (0.608683577,1.29468583) (0.615228561,1.29947625) (0.621773546,1.3042539) (0.628318531,1.30901699) (0.634863515,1.31376372) (0.6414085,1.31849228) (0.647953485,1.32320086) (0.654498469,1.32788764) (0.661043454,1.33255079) (0.667588439,1.3371885) (0.674133424,1.34179892) (0.680678408,1.34638023) (0.687223393,1.35093058) (0.693768378,1.35544814) (0.700313362,1.35993105) (0.706858347,1.36437747) (0.713403332,1.36878554) (0.719948316,1.37315342) (0.726493301,1.37747926) (0.733038286,1.38176119) (0.739583271,1.38599736) (0.746128255,1.39018592) (0.75267324,1.39432502) (0.759218225,1.39841279) (0.765763209,1.40244738) (0.772308194,1.40642694) (0.778853179,1.41034962) (0.785398163,1.41421356) (0.791943148,1.41801693) (0.798488133,1.42175787) (0.805033117,1.42543454) (0.811578102,1.42904511) (0.818123087,1.43258773) (0.824668072,1.43606059) (0.831213056,1.43946184) (0.837758041,1.44278968) (0.844303026,1.44604227) (0.85084801,1.44921782) (0.857392995,1.45231452) (0.86393798,1.45533056) (0.870482964,1.45826416) (0.877027949,1.46111353) (0.883572934,1.46387689) (0.890117919,1.46655247) (0.896662903,1.46913851) (0.903207888,1.47163326) (0.909752873,1.47403497) (0.916297857,1.4763419) (0.922842842,1.47855233) (0.929387827,1.48066455) (0.935932811,1.48267683) (0.942477796,1.4845875) (0.949022781,1.48639486) (0.955567765,1.48809724) (0.96211275,1.48969298) (0.968657735,1.49118042) (0.97520272,1.49255794) (0.981747704,1.4938239) (0.988292689,1.49497669) (0.994837674,1.49601471) (1.00138266,1.49693638) (1.00792764,1.49774013) (1.01447263,1.49842439) (1.02101761,1.49898763) (1.0275626,1.49942831) (1.03410758,1.49974493) (1.04065257,1.49993599) (1.04719755,1.5) (1.05374254,1.4999355) (1.06028752,1.49974105) (1.06683251,1.4994152) (1.07337749,1.49895656) (1.07992247,1.49836371) (1.08646746,1.49763528) (1.09301244,1.49676991) (1.09955743,1.49576625) (1.10610241,1.49462299) (1.1126474,1.49333881) (1.11919238,1.49191243) (1.12573737,1.49034258) (1.13228235,1.48862802) (1.13882734,1.48676752) (1.14537232,1.48475987) (1.15191731,1.48260389) (1.15846229,1.48029841) (1.16500728,1.47784229) (1.17155226,1.4752344) (1.17809725,1.47247365) (1.18464223,1.46955895) (1.19118721,1.46648925) (1.1977322,1.46326351) (1.20427718,1.45988072) (1.21082217,1.4563399) (1.21736715,1.45264008) (1.22391214,1.44878031) (1.23045712,1.44475968) (1.23700211,1.44057729) (1.24354709,1.43623227) (1.25009208,1.43172378) (1.25663706,1.42705098) (1.26318205,1.42221309) (1.26972703,1.41720934) (1.27627202,1.41203897) (1.282817,1.40670126) (1.28936198,1.40119551) (1.29590697,1.39552106) (1.30245195,1.38967726) (1.30899694,1.38366349) (1.31554192,1.37747916) (1.32208691,1.3711237) (1.32863189,1.36459657) (1.33517688,1.35789725) (1.34172186,1.35102527) (1.34826685,1.34398015) (1.35481183,1.33676148) (1.36135682,1.32936884) (1.3679018,1.32180185) (1.37444679,1.31406018) (1.38099177,1.30614348) (1.38753676,1.29805148) (1.39408174,1.2897839) (1.40062672,1.2813405) (1.40717171,1.27272108) (1.41371669,1.26392545) (1.42026168,1.25495346) (1.42680666,1.24580498) (1.43335165,1.23647992) (1.43989663,1.22697821) (1.44644162,1.21729981) (1.4529866,1.20744472) (1.45953159,1.19741294) (1.46607657,1.18720453) (1.47262156,1.17681956) (1.47916654,1.16625814) (1.48571153,1.15552041) (1.49225651,1.14460653) (1.4988015,1.13351669) (1.50534648,1.12225112) (1.51189146,1.11081006) (1.51843645,1.09919381) (1.52498143,1.08740266) (1.53152642,1.07543697) (1.5380714,1.06329709) (1.54461639,1.05098343) (1.55116137,1.03849642) (1.55770636,1.02583652) (1.56425134,1.0130042) (1.57079633,1)};
\addlegendentry{量子预测}
\addplot+[mark=none] coordinates{(0,1) (0.00654498469,1) (0.0130899694,1) (0.0196349541,1) (0.0261799388,1) (0.0327249235,1) (0.0392699082,1) (0.0458148929,1) (0.0523598776,1) (0.0589048623,1) (0.0654498469,1) (0.0719948316,1) (0.0785398163,1) (0.085084801,1) (0.0916297857,1) (0.0981747704,1) (0.104719755,1) (0.11126474,1) (0.117809725,1) (0.124354709,1) (0.130899694,1) (0.137444679,1) (0.143989663,1) (0.150534648,1) (0.157079633,1) (0.163624617,1) (0.170169602,1) (0.176714587,1) (0.183259571,1) (0.189804556,1) (0.196349541,1) (0.202894526,1) (0.20943951,1) (0.215984495,1) (0.22252948,1) (0.229074464,1) (0.235619449,1) (0.242164434,1) (0.248709418,1) (0.255254403,1) (0.261799388,1) (0.268344372,1) (0.274889357,1) (0.281434342,1) (0.287979327,1) (0.294524311,1) (0.301069296,1) (0.307614281,1) (0.314159265,1) (0.32070425,1) (0.327249235,1) (0.333794219,1) (0.340339204,1) (0.346884189,1) (0.353429174,1) (0.359974158,1) (0.366519143,1) (0.373064128,1) (0.379609112,1) (0.386154097,1) (0.392699082,1) (0.399244066,1) (0.405789051,1) (0.412334036,1) (0.41887902,1) (0.425424005,1) (0.43196899,1) (0.438513975,1) (0.445058959,1) (0.451603944,1) (0.458148929,1) (0.464693913,1) (0.471238898,1) (0.477783883,1) (0.484328867,1) (0.490873852,1) (0.497418837,1) (0.503963822,1) (0.510508806,1) (0.517053791,1) (0.523598776,1) (0.53014376,1) (0.536688745,1) (0.54323373,1) (0.549778714,1) (0.556323699,1) (0.562868684,1) (0.569413668,1) (0.575958653,1) (0.582503638,1) (0.589048623,1) (0.595593607,1) (0.602138592,1) (0.608683577,1) (0.615228561,1) (0.621773546,1) (0.628318531,1) (0.634863515,1) (0.6414085,1) (0.647953485,1) (0.654498469,1) (0.661043454,1) (0.667588439,1) (0.674133424,1) (0.680678408,1) (0.687223393,1) (0.693768378,1) (0.700313362,1) (0.706858347,1) (0.713403332,1) (0.719948316,1) (0.726493301,1) (0.733038286,1) (0.739583271,1) (0.746128255,1) (0.75267324,1) (0.759218225,1) (0.765763209,1) (0.772308194,1) (0.778853179,1) (0.785398163,1) (0.791943148,1) (0.798488133,1) (0.805033117,1) (0.811578102,1) (0.818123087,1) (0.824668072,1) (0.831213056,1) (0.837758041,1) (0.844303026,1) (0.85084801,1) (0.857392995,1) (0.86393798,1) (0.870482964,1) (0.877027949,1) (0.883572934,1) (0.890117919,1) (0.896662903,1) (0.903207888,1) (0.909752873,1) (0.916297857,1) (0.922842842,1) (0.929387827,1) (0.935932811,1) (0.942477796,1) (0.949022781,1) (0.955567765,1) (0.96211275,1) (0.968657735,1) (0.97520272,1) (0.981747704,1) (0.988292689,1) (0.994837674,1) (1.00138266,1) (1.00792764,1) (1.01447263,1) (1.02101761,1) (1.0275626,1) (1.03410758,1) (1.04065257,1) (1.04719755,1) (1.05374254,1) (1.06028752,1) (1.06683251,1) (1.07337749,1) (1.07992247,1) (1.08646746,1) (1.09301244,1) (1.09955743,1) (1.10610241,1) (1.1126474,1) (1.11919238,1) (1.12573737,1) (1.13228235,1) (1.13882734,1) (1.14537232,1) (1.15191731,1) (1.15846229,1) (1.16500728,1) (1.17155226,1) (1.17809725,1) (1.18464223,1) (1.19118721,1) (1.1977322,1) (1.20427718,1) (1.21082217,1) (1.21736715,1) (1.22391214,1) (1.23045712,1) (1.23700211,1) (1.24354709,1) (1.25009208,1) (1.25663706,1) (1.26318205,1) (1.26972703,1) (1.27627202,1) (1.282817,1) (1.28936198,1) (1.29590697,1) (1.30245195,1) (1.30899694,1) (1.31554192,1) (1.32208691,1) (1.32863189,1) (1.33517688,1) (1.34172186,1) (1.34826685,1) (1.35481183,1) (1.36135682,1) (1.3679018,1) (1.37444679,1) (1.38099177,1) (1.38753676,1) (1.39408174,1) (1.40062672,1) (1.40717171,1) (1.41371669,1) (1.42026168,1) (1.42680666,1) (1.43335165,1) (1.43989663,1) (1.44644162,1) (1.4529866,1) (1.45953159,1) (1.46607657,1) (1.47262156,1) (1.47916654,1) (1.48571153,1) (1.49225651,1) (1.4988015,1) (1.50534648,1) (1.51189146,1) (1.51843645,1) (1.52498143,1) (1.53152642,1) (1.5380714,1) (1.54461639,1) (1.55116137,1) (1.55770636,1) (1.56425134,1) (1.57079633,1)};
\addlegendentry{局域界}
\end{axis}\end{tikzpicture}\end{center}

\end{solution}
\originalsection{第2题：CHSH游戏}
\begin{exercise}
Alice,Bob共享$\ket{\Phi^+}=(\ket{++}+\ket{--})/\sqrt2$, 独立收到均匀比特$x,y$, 输出$a,b$, 要求$a\oplus b=xy$. (a)证测量基$\ket{v_a},\ket{w_b}$下$P(a,b)=|\braket{\bar v_a}{w_b}|^2/2$. (b)证明经典策略最多赢3/4. (c)实平面基由转角$\alpha_x,\beta_y$决定, 求同输出概率. (d)选角实现量子更高胜率.
\end{exercise}
\begin{solution}
(a)两bra作用Bell态得$(v_{a+}^*w_{b+}^*+v_{a-}^*w_{b-}^*)/\sqrt2$, 模平方等于$|\braket{\bar v_a}{w_b}|^2/2$, 横线对系数复共轭. (b)每人一个比特输入的确定策略有$0,1,x,1-x$四种. 若四输入均赢, 把四个等式在模2相加, 左侧每个输出出现两次得0, 右侧$xy$真值表和1, 矛盾. 常输出0在前三输入赢, 达3/4. 共享随机只是确定策略的凸组合, 不能超过最大值.

(c)取输出0单位矢$(\cos\frac\alpha2,\sin\frac\alpha2)$, 输出1其正交矢, 两人的夹角差为半角, 相等输出总概率$\cos^2[(\alpha_x-\beta_y)/2]$. (d)胜率为
\[
P_{\rm win}=\frac14\left[\cos^2\frac{\alpha_0-\beta_0}2+
\cos^2\frac{\alpha_0-\beta_1}2+
\cos^2\frac{\alpha_1-\beta_0}2+
\sin^2\frac{\alpha_1-\beta_1}2\right].
\]
选$\alpha_0=0,\alpha_1=\pi/2,\beta_0=\pi/4,\beta_1=-\pi/4$, 四项都$\cos^2(\pi/8)$, 胜率$(2+\sqrt2)/4\simeq0.853553$. 各地输出边缘分布仍均匀, 相关增强不能用来超光速传信.
\end{solution}
\originalsection{第3题：MIT角动量讲义练习1--7}
\begin{exercise}
按官方Chapter 9的Exercise 1--7, 依次证明向量积的伴随、有序叉积平方、$p\times L$交换式、其Hermitian部分、固定轴生成旋转、向量点积为标量/叉积仍为向量、$L\times L=\ii\hbar L$.
\end{exercise}
\begin{solution}
\textbf{练习1.} 分量$(a\cdot b)^\dagger=\sum b_i^\dagger a_i^\dagger=b^\dagger\cdot a^\dagger$. $(a\times b)_i=\epsilon_{ijk}a_jb_k$, 伴随后交换哑指标$j,k$引负号, 得$(a\times b)^\dagger=-b^\dagger\times a^\dagger$. 次序是必需的.

\textbf{练习2.} 两个$\epsilon$缩并给$\sum_{jk}(a_jb_ka_jb_k-a_jb_ka_kb_j)$. 第一项将$b_k$移过$a_j$得到$a^2b^2-a_j[a_j,b_k]b_k$. 第二项先把$b_k$移过$a_k$, 再把$b_k$与$b_j$换序并把$a_k$移过$b_j$, 整理为
\begin{align*}
(a\times b)^2={}&a^2b^2-(a\cdot b)^2-a_j[a_j,b_k]b_k
+a_j[a_k,b_k]b_j\\
&-a_j[a_k,b_j]b_k-a_ja_k[b_k,b_j].
\end{align*}
重复指标求和. 若$[a_i,b_j]=\gamma\delta_{ij}$、$[b_i,b_j]=0$, 三个剩余项系数为$-\gamma+3\gamma-\gamma=\gamma$, 得$a^2b^2-(a\cdot b)^2+\gamma a\cdot b$. 对$r,p$,$\gamma=\ii\hbar$, 即$L^2=r^2p^2-(r\cdot p)^2+\ii\hbar r\cdot p$.

\textbf{练习3.} $[L_i,p_j]=\ii\hbar\epsilon_{ijk}p_k$, 所以$(p\times L+L\times p)_i=\epsilon_{ijk}[p_j,L_k]=2\ii\hbar p_i$. 即$p\times L=-L\times p+2\ii\hbar p$. \textbf{练习4.} 用练习1,$(p\times L)^\dagger=-L\times p$, 故Hermitian部分$(p\times L-L\times p)/2=p\times L-\ii\hbar p$.

\textbf{练习5.} $[n\cdot L,u_j]=n_i\ii\hbar\epsilon_{ijk}u_k=-\ii\hbar(n\times u)_j$, 对任意向量算符成立. \textbf{练习6.} 乘积对易规则使$[L_i,u_jv_j]=\ii\hbar\epsilon_{ijk}(u_kv_j+u_jv_k)=0$, 两项交换$j,k$抵消且没有交换算符自身顺序. 对叉积,$[L_i,\epsilon_{jmn}u_mv_n]=\ii\hbar\epsilon_{jmn}(\epsilon_{imk}u_kv_n+\epsilon_{ink}u_mv_k)$, 缩并为$\ii\hbar\epsilon_{ijk}(u\times v)_k$, 故仍向量.

\textbf{练习7.} $(L\times L)_i=\epsilon_{ijk}[L_j,L_k]/2=\ii\hbar\epsilon_{ijk}\epsilon_{jkn}L_n/2=\ii\hbar L_i$. 反向$\epsilon$缩并也恢复全部角动量对易代数, 两写法等价. 七题均有MIT实际题面, 区别于本篇第4题的外书仅题号.
\end{solution}
\originalsection{第5题：二维谐振子的隐藏SU(2)}
\begin{exercise}
二维各向同性振子的$a_L=(a_x+\ii a_y)/\sqrt2$, $a_R=(a_x-\ii a_y)/\sqrt2$, 物理轨道$L_z=\hbar(N_R-N_L)$. 独立圆模式$R,L$的$H=\hbar\omega(N_R+N_L+1)$, 定义$J_z=\alpha\hbar(N_R-N_L)$,$J_+=\beta\hbar a_R^\dagger a_L$,$J_-=\beta\hbar a_L^\dagger a_R$, $\alpha,\beta$实. 求标准角动量代数的参数, 证与$H$对易并分类总激发$n$子空间的$j,m$与简并. 为什么不能取$\alpha=1$?
\end{exercise}
\begin{solution}
$[N_R-N_L,a_R^\dagger a_L]=2a_R^\dagger a_L$给$[J_z,J_+]=2\alpha\hbar J_+$, 标准要求$\alpha=1/2$. $[a_R^\dagger a_L,a_L^\dagger a_R]=N_R-N_L$给$\beta^2=1$, 选$\beta=1$, $-1$为共同相位改变. 这些操作只在两模式间搬一个量子, 保持总$n$, 因而均与$H$对易.

$(a_R^\dagger)^n\ket0$是$m=n/2$的最高态, $J_+$杀死它. 从$J^2=J_-J_++J_z^2+\hbar J_z$得$j=n/2$, 再降$n$次穷尽$n+1$个态. 一般$m=(n_R-n_L)/2$, $n_R=n/2+m,n_L=n/2-m$, $m=-j,\ldots,j$, 能量$(n+1)\hbar\omega$. 若$\alpha=1$, 搬一个量子使$m$变2而非标准单位1, 不满足角动量代数; 也会把最高$m$误成$n$却只有$n+1$态, 与$2j+1$不符.
\end{solution}
